\section{Further results}
\label{app:results}

This appendix gives supporting analyses for Section~\ref{sec:results}. Every number is a macro written from stored result
files.

\subsection{Structural identifiability and rank}
\label{apx:structural}\label{app:structural}

StructuralIdentifiability.jl, in versions \texttt{\juliaversion} and \texttt{\siversion} and with glucose as the output, asks
what an ideal observation of glucose determines. Seven formulations were run; the four of Table~\ref{tab:structural} differ in
outcome, and the raw output of every run is archived with the code. The
answer depends on what is assumed about the gut at the start of a meal, an assumption the data cannot settle
(Table~\ref{tab:structural}). With a generic, known intestinal content the tool reports $\ke$ and $\ka$ globally
identifiable, which is the formulation fixed in the plan. With the gut state unconstrained or unknown they are only
locally identifiable. In canonical coordinates the symmetric functions $\ke+\ka$ and $\ke\ka$ are globally
identifiable, which is the unordered pair.

\begin{table}[tbp]
\centering\footnotesize
\begin{tabularx}{\linewidth}{Xccc}
\toprule
formulation & initial gut state & $\SI$ & $\ke,\ka$ \\
\midrule
meal as input, all initial conditions generic and known & known, generic & global & global \\
\emph{planned}: meal as initial stomach content, no input & known, generic & global & global \\
meal as input & unconstrained or unknown & global & local \\
canonical coordinates $(\ke+\ka,\ \ke\ka)$ & unconstrained or unknown & global & the two combinations are global \\
\bottomrule
\end{tabularx}
\caption{Structural identifiability of the model with glucose as output, by formulation. A generic known
intestinal content makes the rates globally identifiable; an unconstrained one leaves the unordered pair.}
\label{tab:structural}
\end{table}

The analysis plan said to drop the exchange framing if the rates were globally identifiable. That condition holds for
the formulation written in the plan and not for the one the simulator runs, so we record the test as inconclusive. We
keep the framing because the empty-gut model that the simulator runs is exactly symmetric, and free-living data
approximate that model to the extent that meals are isolated.

The stacked sensitivity matrix has generic rank three at relative tolerance $10^{-6}$ in \resRankFullIauc{} of
\resRankDrawsIauc{} random draws for iAUC, and in every draw for iAUC with the centroid and for the trace. All three
parameters are locally identifiable under every target of the ladder. What differs is conditioning: the ratio of the
smallest to the largest singular value is \resRankRatioIauc{} for iAUC, \resRankRatioCentroid{} with the centroid and
\resRankRatioTrace{} for the trace. We call iAUC locally identifiable but ill-conditioned and claim nothing stronger
than the exchange symmetry.

\subsection{Timing under the area: estimates, Fisher information, and cohorts}
\label{apx:timing}

\paragraph{Estimates on bounds.}
At the primary box the estimates of $\ke$ and $\ka$ are on a bound in \resPinnedFracKeOne\% and \resPinnedFracKaOne\%
of the \resSubjectsOne{} CGMacros subjects, and on the upper bound in \resPinnedUpperKeOne{} and
\resPinnedUpperKaOne{}. For $\SI$ the figure is \resPinnedFracSIOne\%. Doubling the box lowers the fractions for $\ke$
and $\ka$ to \resPinnedFracKeTwo\% and \resPinnedFracKaTwo\%, and halving it raises them to \resPinnedFracKeHalf\% and
\resPinnedFracKaHalf\%. They never reach zero (Figure~\ref{fig:gutsweep}b). The negative log-likelihood of the
first-stage estimate exceeds that of the best of six L-BFGS runs by at most
\resOptGapMaxHalf, \resOptGapMaxOne{} and \resOptGapMaxTwo{} at the three widths, always below $\delta$.

\paragraph{Fisher information.}
In log coordinates the median eigenvalues of the information matrix at the primary box are \resEigMedianOne,
\resEigMedianTwo{} and \resEigMedianThree{} (Figure~\ref{fig:fisher}a). The weakest eigenvector of the timing block
lies along the exchange direction, with a median absolute cosine of \resCosineIaucOne{} under iAUC,
\resCosineCentroidOne{} with the centroid and \resCosineTraceOne{} under the trace.

\paragraph{Pinning is partly an area deficit.}
Subjects whose $\ke$ or $\ka$ sits on the upper bound are not simply those with a flat likelihood. Their mean observed
iAUC is \resBiasObsPinnedKe{} against \resBiasObsOtherKe{} for the others, and the model under-predicts them by
\resBiasSignedKe{} mg\,dL$^{-1}$\,min on average, against \resBiasSignedOtherKe{} for the others. The model cannot
deliver the observed area at any admissible timing, so the fit pushes the rates to their limit. So part of the pinning reflects misfit of the model or of the carbohydrate logs.

\paragraph{Three cohorts.}
The Shanghai cohort (\resRepShanghaiOneN{} subjects, $15$-minute sensor, dietary records) and the cohort with
standardized meals (\resRepHallOneN{} subjects) show the same picture (Table~\ref{tab:cohortresults}). No interval for $\ke$ or $\ka$ is bounded in either cohort at the primary or the double box.
The weakest Fisher direction has $|\cos|\ge0.8$ with the exchange direction in \resRepShanghaiOneCosine\% of Shanghai subjects and \resRepHallOneCosine\% of standardized-meal subjects. Most estimates in these cohorts sit on a bound at the
primary box, so only \resRepShanghaiOneInteriorKe{} and \resRepShanghaiOneInteriorKa{} Shanghai subjects and
\resRepHallOneInteriorKe{} and \resRepHallOneInteriorKa{} subjects of the other cohort are interior for $\ke$ and
$\ka$. The replication rests on the Fisher alignment and the flat profiles, not on the interior fractions.

\begin{table}[t]
\centering\footnotesize\setlength{\tabcolsep}{4pt}
\begin{tabular}{llrcccc}
\toprule
cohort & box & subjects & estimate on a bound (\%) & bounded interval (\%) & interior for $\ke$ / $\ka$ & $|\cos|\ge0.8$ (\%) \\
 & & & $\SI$ / $\ke$ / $\ka$ & $\SI$ / $\ke$ / $\ka$ & & \\
\midrule
CGMacros & $1\times$ & 45 & 51 / 89 / 87 & 7 / 0 / 0 & 5 / 6 & 100 \\
CGMacros & $2\times$ & 45 & 11 / 58 / 76 & 53 / 0 / 0 & 19 / 11 & 100 \\
Shanghai & $1\times$ & 87 & 45 / 95 / 97 & 7 / 0 / 0 & 4 / 2 & 100 \\
Shanghai & $2\times$ & 87 & 10 / 68 / 72 & 45 / 0 / 0 & 24 / 24 & 100 \\
Standardized meals & $1\times$ & 15 & 73 / 80 / 67 & 0 / 0 / 0 & 3 / 5 & 100 \\
Standardized meals & $2\times$ & 15 & 7 / 60 / 40 & 0 / 0 / 0 & 6 / 9 & 100 \\
\bottomrule
\end{tabular}
\caption{Replication of the identifiability result on three cohorts, iAUC objective. The last column is the share of subjects whose weakest timing Fisher direction has an absolute cosine of at least $0.8$ with the exchange direction. Interior counts are subjects whose estimate of the parameter is more than one percent of the box width from both bounds.}
\label{tab:cohortresults}
\end{table}


\paragraph{Why the box matters.}
Most $\SI$ intervals are one-sided, so the data constrain a direction, and the box decides whether the other side is
bounded. Part of this dependence is saturation of the response. The slope of the iAUC with respect to $\SI$ at the top
of the primary box is \resSatTopToBottom{} of its value at the bottom, and at the top of the double box it is
\resSatDoubleTop{} of that at the top of the primary box (exploratory). Because $\gamma$, $p_3$, and $\SI$ enter only
through their product, rescaling either constant in the generator rescales the fitted $\SI$ by the same factor.

\paragraph{Spacing of logged meals.}\label{apx:phase13}
The CGMacros logs contain \resOverlapGaps{} gaps between consecutive meals of a subject, with a median of
\resOverlapMedianGap{} hours. The share under three, four and six hours is \resOverlapWithinThree\% (Wilson interval
\resOverlapWithinThreeLow\ to \resOverlapWithinThreeHigh\%), \resOverlapWithinFour\% (\resOverlapWithinFourLow\ to
\resOverlapWithinFourHigh\%) and \resOverlapWithinSix\% (\resOverlapWithinSixLow\ to \resOverlapWithinSixHigh\%).
The model assumes an empty gut at each meal. We did not measure the residual content and so do not say how much of
the exchange symmetry this breaks.

\begin{figure}[tbp]
\centering
\safefigure[width=0.80\linewidth]{a1_gut_and_boxes}
\caption{\textbf{(a)} Largest and median change in simulated glucose under the exchange $\ke\leftrightarrow\ka$
as a function of residual intestinal content, as a percentage of the meal (shaded: $5$ to $10$\,mg\,dL$^{-1}$, for scale).
\textbf{(b)} Subjects whose estimate lies on a bound, by box width.}
\label{fig:gutsweep}
\end{figure}

\subsection{Held-out prediction: further detail}
\label{apx:prediction}

Table~\ref{tab:prediction} gives the held-out iAUC error of every cell and Figure~\ref{fig:prediction} the paired
differences. With $\SI$ only, gradient fitting and grid search differ by \resDiffGradOneVsGridOne{}
(\resDiffGradOneVsGridOneLow{} to \resDiffGradOneVsGridOneHigh). The training loss of the gradient fit is within one
percent of the grid minimum in \resGapWithinTolerance\% of folds, and a budget of $500$ steps changes no
verdict (difference \resStepsGradThreeFiveHundredVsGridThree{}, interval
\resStepsGradThreeFiveHundredVsGridThreeLow{} to \resStepsGradThreeFiveHundredVsGridThreeHigh). The trace objective is
\resDiffGradThreeTraceVsGradThree{} (\resDiffGradThreeTraceVsGradThreeLow{} to \resDiffGradThreeTraceVsGradThreeHigh)
worse than the iAUC objective on iAUC error. It improves what it targets: the held-out trace error falls from
\resTraceRmseGradThree{} to \resTraceRmseGradThreeTrace{} mg\,dL$^{-1}$ and the peak-time error from
\resPeakGradThree{} to \resPeakGradThreeTrace{} min. The personal mean, which uses no model, reaches
\resTraceRmsePersonalMean{} and \resPeakPersonalMean{}, so the model adds little there either.

Against the personal mean, the gradient fit has lower iAUC error (\resDiffGradThreeVsPersonalMean{},
\resDiffGradThreeVsPersonalMeanLow{} to \resDiffGradThreeVsPersonalMeanHigh, better in
\resWinsGradThreeVsPersonalMean{} of \resSubjectsOne{} subjects, Holm-adjusted $p=\resPHolmGradThreeVsPersonalMean$).

\begin{figure}[tbp]
\centering
\safefigure[width=0.78\linewidth]{f10_prediction_vs_parameters}
\caption{Held-out iAUC error of the three-parameter fit divided by the error of the model with the true parameters under the same noise and carbohydrate error (the true-parameter reference, which is not a lower bound),
on the true-model replica with fitted truths (seed~$0$) and with random truths (seeds~$0$ and~$1$).
\textbf{(a)} Per subject, against the root mean squared log error of $\ke$ and $\ka$ in the iAUC estimate. A model
whose error rose with its parameter error would slope upward. \textbf{(b)} Distribution over subjects; diamonds:
median error of the personal mean relative to the same reference.}
\label{fig:pvp}
\end{figure}

\paragraph{Shanghai cohort.}
On Shanghai (\resShCvSubjects{} subjects, \resShCvRepeats{} repeats), the three-parameter fit has lower held-out iAUC error than the one-parameter fit in \resShCvWinsGradThreeVsGradOne{} of \resShCvSubjects{} subjects (by \resShCvDiffGradThreeVsGradOne{}, interval excluding zero), and both timing-parameter comparisons (\resShCvDiffGradThreeVsGridThree{} for gradient vs.\ grid and \resShCvDiffGradThreeVsGradOne{} for three parameters vs.\ $\SI$ alone) lie inside the $\pm 150$ equivalence margin (\resShCvMarginShare\% of the personal-mean error of \resShCvMaePersonalMean{}).
Against the personal mean, the paired difference is \resShCvDiffGradThreeVsPersonalMean{} (\resShCvDiffGradThreeVsPersonalMeanLow{} to \resShCvDiffGradThreeVsPersonalMeanHigh), an interval that includes both zero and differences exceeding the margin, so these data do not settle whether the simulator beats a model-free baseline on Shanghai.

\paragraph{Chronological folds.}
Folds that respect time give the same comparison. The three-parameter fit differs from the personal mean by
\resChronoRandomDiffMean{} mg\,dL$^{-1}$\,min with random folds, by \resChronoBlockedDiffMean{}
(\resChronoBlockedDiffMeanLow{} to \resChronoBlockedDiffMeanHigh) with contiguous blocks of meals held out, and by
\resChronoForwardDiffMean{} (\resChronoForwardDiffMeanLow{} to \resChronoForwardDiffMeanHigh) when only earlier meals train.

\begin{table}[t]
\centering\small
\begin{tabular}{lcccc}
\toprule
cell & iAUC, CGMacros & iAUC, Shanghai & peak time, CGMacros & trace, CGMacros \\
 & (mg\,dL$^{-1}$\,min) & (mg\,dL$^{-1}$\,min) & (min) & (mg\,dL$^{-1}$) \\
\midrule
gradient, three parameters & 3025 & 3854 & 47.5 & 35.9 \\
grid search, three parameters & 3026 & 3855 & 48.0 & 36.3 \\
gradient, $\SI$ only & 3059 & 3908 & 44.4 & 33.1 \\
grid search, $\SI$ only & 3059 & 3908 & 44.4 & 33.1 \\
gradient, trace objective & 3057 & -- & 41.6 & 31.2 \\
personal mean & 3257 & 3892 & 41.0 & 32.0 \\
population default & 3359 & 4220 & 46.8 & 34.8 \\
persistence & 4507 & 5535 & 54.6 & 46.5 \\
\bottomrule
\end{tabular}
\caption{Held-out mean absolute error (iAUC, peak time) and root mean squared error (trace) by cell; mean over subjects of the mean over five repeats of five-fold cross-validation. The Shanghai column is empty where the cell was not run.}
\label{tab:prediction}
\end{table}


\begin{figure}[tbp]
\centering
\safefigure[width=0.80\linewidth]{f7_prediction}
\caption{Paired differences in held-out iAUC error between cells (first minus second; negative favors the first), with
$95\%$ bootstrap intervals over subjects. Shaded: the equivalence margin of $150$. Blue: CGMacros; green: Shanghai.}
\label{fig:prediction}
\end{figure}

\paragraph{What the remaining error is.}
We do not report additive shares of the held-out error for noise, optimization, and target, because no definition of the
decomposition makes the terms add. We report what-if noise budgets on the replica instead (Figure~\ref{fig:replica}c).
The ratio of the replica's held-out iAUC error to the real data's is \resReplicaRatioCgmOffCarbZero{} with no noise,
\resReplicaRatioCgmOffCarbQuarter{} with $25\%$ carbohydrate error alone, \resReplicaRatioCgmOnCarbZero{} with CGM noise
alone, \resReplicaRatioCgmOnCarbQuarter{} (\resReplicaRatioCgmOnCarbQuarterLow--\resReplicaRatioCgmOnCarbQuarterHigh)
with both, and \resReplicaRatioCgmOnCarbHalf{} with $50\%$ carbohydrate error. Grid search and gradient fitting give the same
error on the replica, so the optimizer is not the source. The CGM noise is estimated from the residuals of the real
fit, which contain model misfit, so these settings are an upper bound on what sensor noise alone would explain.

\subsection{Targets, window, and the combination of the rates}
\label{apx:more}

\paragraph{The window.}
The sensitivity of the iAUC to $\ke$ and $\ka$, relative to its sensitivity to $\SI$, depends on the window
(Figure~\ref{fig:window}). At ninety minutes the area is more sensitive to timing than to $\SI$, with median ratios of
$\resWinNinetyLinKe$ and $\resWinNinetyLinKa$ for $\ke$ and $\ka$ in the linearized model. At the three-hour window
used here the ratios are $\resWinOneEightyLinKe$ and $\resWinOneEightyLinKa$, and at six hours
$\resWinThreeSixtyLinKe$ and $\resWinThreeSixtyLinKa$. The nonlinear engine agrees at three
hours ($\resWinOneEightyNlKe$, $\resWinOneEightyNlKa$) but shows slight non-monotonicity at $6$\,h due to the softplus baseline floor. A shorter window increases timing sensitivity relative to $\SI$ while reducing $\SI$ sensitivity. Testing this directly in profile likelihood across $T\in\{60, 90, 120, 180\}$\,min on CGMacros (primary box, earlier $15$-point protocol; a rerun under the converged protocol is pending) shows that shortening the window moves gut-rate estimates off the upper bound (\resWinProfSixtyKeInt{} and \resWinProfSixtyKaInt{} interior for $\ke$ and $\ka$ at $60$\,min; \resWinProfNinetyKeInt{} and \resWinProfNinetyKaInt{} at $90$\,min; \resWinProfOneTwentyKeInt{} and \resWinProfOneTwentyKaInt{} at $120$\,min; $5$ and $6$ at $180$\,min), yet \resWinProfSixtyKeBd{} subjects have a bounded interval for $\ke$ or $\ka$ at any window, while bounded $\SI$ intervals drop from \resProfileIaucSIInteriorBoundedOne{} of \resProfileIaucSIInteriorNOne{} at $180$\,min to \resWinProfOneTwentySIBd{} of \resWinProfOneTwentySIInt{} ($120$\,min), \resWinProfNinetySIBd{} of \resWinProfNinetySIInt{} ($90$\,min), and \resWinProfSixtySIBd{} of \resWinProfSixtySIInt{} ($60$\,min).

\begin{figure}[tbp]
\centering
\safefigure[width=0.78\linewidth]{f8_window}
\caption{Median (line) and interquartile range (band) over subjects of
$|\partial\,\iauc/\partial\log k|/|\partial\,\iauc/\partial\log\SI|$ for $k=\ke$ and $k=\ka$, against the
window after the meal, for the linearized and the nonlinear engine. Dashed: the thresholds of H2; dotted:
a ratio of one.}
\label{fig:window}
\end{figure}

\paragraph{The ladder of targets.}
Information about timing, measured by the condition number of the Schur-complement block after profiling out $\SI$,
rises with the richness of the target at a common reference point: \resCondRefIaucOne{} for iAUC,
\resCondRefCentroidOne{} with the centroid, and \resCondRefTraceOne{} for the trace at the primary box
(Figure~\ref{fig:fisher}b). It does not rise at each target's own estimate (\resCondOwnIaucOne{},
\resCondOwnCentroidOne{}, \resCondOwnTraceOne{}), where those fits sit on bounds and the information is nearly singular.

\begin{figure}[tbp]
\centering
\safefigure[width=0.78\linewidth]{f2_fisher}
\caption{\textbf{(a)} Eigenvalues of the Fisher information in log coordinates evaluated at each target's own fitted estimate, by observable, primary
box. \textbf{(b)} Median condition number of the timing block evaluated at a common reference point across targets, by observable
and box.}
\label{fig:fisher}
\end{figure}

\paragraph{The tied-rate model.}
A model with $\ke=\ka$, which has only $\tau_1$ and $\SI$, predicts as well as the three-parameter model (difference in
held-out iAUC error \resTiedDiff{} mg\,dL$^{-1}$\,min, interval \resTiedDiffLow{} to \resTiedDiffHigh). The coordinate
fits sit on the tied boundary $p=\tau_1^2/4$ in \resBoundaryCoords\% of cross-validation folds under iAUC and
\resBoundaryCoordsTrace\% under the trace, so for many subjects the data are consistent with equal rates.

\subsection{Numerical sensitivity of the profiles}
\label{apx:numerics}

The trace profiles of Section~\ref{sec:more} use a $29$-point grid with every point solved to convergence; a point
where the L-BFGS-B line search fails is solved again from two other starting points and the lower value is kept.
The checks below were run on an earlier protocol, a $15$-point grid with a fixed-step inner solver (Adam, $100$ steps),
and show why we left it.

\paragraph{Inner budget.}
Table~\ref{tab:inner} compares the profiles at $100$ and $500$ inner steps (same estimates, same grids). A grid point is flagged
when its inner optimization was still improving by more than $0.05$ in negative log-likelihood. The longer run settles nearly
every point and changes no profile class (\resInnerQuarterChanged{} of \resInnerQuarterChangedOf{} at a quarter of the noise,
\resInnerRealChanged{} of \resInnerRealChangedOf{} on the real data), and the real-data bounded fractions stay at
\resInnerRealFiveHundredBoundedTau\% and \resInnerRealFiveHundredBoundedP\%. It does not restore coverage at a quarter of the
noise with exact carbohydrate (\resInnerQuarterFiveHundredCoverTau\% and \resInnerQuarterFiveHundredCoverP\%).

\begin{table}[tbp]
\centering\footnotesize
\begin{tabular}{lrrrrrr}
\toprule
 & \multicolumn{2}{c}{flagged points (\%)} & profile classes & \multicolumn{2}{c}{covers truth, long run (\%)} \\
setting & default & long & changed & $\tau_1$ & $p$ \\
\midrule
real data & \resInnerRealHundredUnsettled & \resInnerRealFiveHundredUnsettled & \resInnerRealChanged{} of \resInnerRealChangedOf & --- & --- \\
replica, error $0.25$ & \resInnerRepHundredUnsettled & \resInnerRepFiveHundredUnsettled & \resInnerRepChanged{} of \resInnerRepChangedOf & \resInnerRepFiveHundredCoverTau & \resInnerRepFiveHundredCoverP \\
replica, exact carbohydrate & \resInnerExactHundredUnsettled & \resInnerExactFiveHundredUnsettled & \resInnerExactChanged{} of \resInnerExactChangedOf & \resInnerExactFiveHundredCoverTau & \resInnerExactFiveHundredCoverP \\
exact, quarter noise & \resInnerQuarterHundredUnsettled & \resInnerQuarterFiveHundredUnsettled & \resInnerQuarterChanged{} of \resInnerQuarterChangedOf & \resInnerQuarterFiveHundredCoverTau & \resInnerQuarterFiveHundredCoverP \\
\bottomrule
\end{tabular}
\caption{Trace coordinate profiles with polished estimates at $100$ and $500$ inner steps, primary box. Profile classes
are bounded, one-sided or flat; each setting has $45$ subjects and two profiles per subject.}
\label{tab:inner}
\end{table}

\paragraph{The likelihood at the true value.}
At a quarter of the noise with exact carbohydrate the likelihood at the true parameters is calibrated: half the
likelihood-ratio statistic has median \resCalQuarterMedianLr{} and ninetieth percentile \resCalQuarterQNinetyLr{}, against
\resCalRefMedian{} and \resCalRefQNinety{} for half a chi-square with three degrees of freedom, and exceeds the cutoff in
\resCalQuarterOver\% of subjects. The stored polished estimate is within the cutoff of the best optimum found from the truth
in \resCalQuarterGapWithin{} of \resCalQuarterN{} subjects, with a largest gap of \resCalQuarterGapMax. With the logging error
the likelihood at the truth is not calibrated (\resCalRepOver\% of subjects above the cutoff).

\paragraph{The profile at the true value.}
Table~\ref{tab:truthprofile} gives the excess of the profile above its reference at the true value of $\tau_1$ and $p$, read
off the stored grid and computed directly there with the other two coordinates solved to convergence. Computed directly, the
excess is within the cutoff for \resTruthQuarterTauDirectWithinCount{} of \resTruthQuarterTauDirectN{} subjects for $\tau_1$
and \resTruthQuarterPDirectWithinCount{} of \resTruthQuarterPDirectN{} for $p$ at a quarter of the noise, whether the solve
starts at the estimate or at the truth and whether it uses Adam or L-BFGS-B; those choices change the median excess by
\resTruthQuarterTauSolver{} (solver) and \resTruthQuarterTauStart{} (start). The excess read off the grid has median
\resTruthQuarterTauStoredMedian{} for $\tau_1$ and \resTruthQuarterPStoredMedian{} for $p$, compared to direct-solve medians of
\resTruthQuarterTauDirectMedian{} for $\tau_1$ and \resTruthQuarterPDirectMedian{} for $p$. The difference is linear
interpolation across a grid that is wide relative to a narrow profile (median \resTruthQuarterTauInterp{} for $\tau_1$ and
\resTruthQuarterPInterp{} for $p$). At the estimated noise, the comparison gives \resTruthExactTauStoredMedian{} vs.\
\resTruthExactTauDirectMedian{} for $\tau_1$, a smaller overstatement that leaves coverage near its nominal level.

\begin{table}[tbp]
\centering\footnotesize
\begin{tabular}{llrrrr}
\toprule
 & & \multicolumn{2}{c}{median excess} & \multicolumn{2}{c}{above the cutoff (\%)} \\
replica & coordinate & grid & direct & grid & direct \\
\midrule
exact, quarter noise & $\tau_1$ & \resTruthQuarterTauStoredMedian & \resTruthQuarterTauDirectMedian & \resTruthQuarterTauStoredOver & \resTruthQuarterTauDirectOver \\
 & $p$ & \resTruthQuarterPStoredMedian & \resTruthQuarterPDirectMedian & \resTruthQuarterPStoredOver & \resTruthQuarterPDirectOver \\
exact carbohydrate & $\tau_1$ & \resTruthExactTauStoredMedian & \resTruthExactTauDirectMedian & \resTruthExactTauStoredOver & \resTruthExactTauDirectOver \\
 & $p$ & \resTruthExactPStoredMedian & \resTruthExactPDirectMedian & \resTruthExactPStoredOver & \resTruthExactPDirectOver \\
error $0.25$ & $\tau_1$ & \resTruthRepTauStoredMedian & \resTruthRepTauDirectMedian & \resTruthRepTauStoredOver & \resTruthRepTauDirectOver \\
 & $p$ & \resTruthRepPStoredMedian & \resTruthRepPDirectMedian & \resTruthRepPStoredOver & \resTruthRepPDirectOver \\
\bottomrule
\end{tabular}
\caption{Excess of the trace coordinate profile at the true value of $\tau_1$ and $p$ (negative log-likelihood units),
interpolated on the stored $15$-point grid and computed directly at the true value. The cutoff is the $1.92$ of the
interval criterion; each row has $45$ subjects.}
\label{tab:truthprofile}
\end{table}

\paragraph{Converged solver and finer grid.}
A converged inner solver (L-BFGS-B at every grid point) changes no profile class (\resConvRealChanged{} of
\resConvRealChangedOf{} on the real data, \resConvQuarterChanged{} of \resConvQuarterChangedOf{} at a quarter of the noise)
and leaves the real-data bounded fractions at \resConvRealTauBounded\% and \resConvRealPBounded\%. A $29$-point grid
in place of $15$ leaves them at \resGridRealTwentyNineTauBounded\% and \resGridRealTwentyNinePBounded\%, changing no
class on the real data (\resGridRealChanged{} of \resGridRealChangedOf{}), the exact-carbohydrate replica (\resGridExactChanged{}),
or the carbohydrate-error replica (\resGridRepChanged{}). Only the narrow profiles of the low-noise replica depend on the grid. On
the $29$-point grid with the converged solver, the nominal interval at a quarter of the noise contains the truth for
\resConvFineTauCover\% of subjects for $\tau_1$ and \resConvFinePCover\% for $p$, the detector bounds $\tau_1$ and $p$ in \resConvFineTauBounded\%
and \resConvFinePBounded\% of interior subjects, and the bounded intervals contain the truth in \resConvFineTauCovbd\% and
\resConvFinePCovbd\%.

\subsection{A second model class}
\label{apx:dm}

The Dalla Man meal model \citep{dallaman2007}, in our implementation, has a three-compartment gut with a nonlinear
gastric-emptying law, so it inherits no exact exchange symmetry. We fitted it by maximum likelihood to the same iAUC data
of the same \resDmN{} subjects with a single-variance Gaussian noise model (exploratory; earlier profile protocol). Its
Fisher information has median effective rank \resDmRankMedian{} of six parameters and condition number
\resDmCondMedian. The profile interval is bounded for $k_{abs}$ in \resDmBoundedKabs\%, for $k_{max}$ in \resDmBoundedKmax\%, and for $k_{min}$ in \resDmBoundedKmin\% of subjects, compared to \resDmBoundedVmx\% (interval
\resDmBoundedVmxLow--\resDmBoundedVmxHigh\%) for its insulin-sensitivity parameter (Figure~\ref{fig:dm}). The estimates
of the three gut parameters lie on a bound of their range in \resDmPinnedKabs, \resDmPinnedKmax, and \resDmPinnedKmin{}
of \resDmN{} subjects, and in the others no interval is bounded. Fitting six parameters (including absorption fraction and sensor lag) to a single scalar summary per meal leaves most directions poorly constrained even without exact exchange symmetry. The claim this
supports is narrow: under the iAUC target, in two model classes that differ in the structure of the gut, the gut
parameters are not recovered from these subjects. It does not say that the exchange explanation carries over, and a
Fisher rank of \resDmRankMedian{} of six shows only that most directions are poorly determined. We ran neither a Dalla Man-on-Dalla-Man replica nor a Dalla Man synthetic study.

\begin{figure}[tbp]
\centering
\safefigure[width=0.80\linewidth]{a3_dalla_man}
\caption{Second model class (Dalla Man), iAUC objective. \textbf{(a)} Eigenvalues of the Fisher information in
log coordinates, over subjects. \textbf{(b)} Subjects with a bounded $95\%$ profile interval per parameter.}
\label{fig:dm}
\end{figure}

\paragraph{A clinical correlation is not a check.}\label{apx:leakage}
In the Shanghai cohort a fitted $\SI$ from an amortized estimator correlated with HbA1c, and most of that correlation
disappeared with fasting glucose partialled out, so we do not use correlations with clinical measures to validate
fitted parameters.
\FloatBarrier
